\documentclass[11pt]{article}
\usepackage{setspace}
\usepackage{geometry}
\geometry{margin=1in}
\setstretch{1.15}

\begin{document}

% --- Applicant Information ---
\noindent
\textbf{Zhongli Wang} \\
Phone: +49 160 9541 6015 \\
Email: nkwangzhongli@gmail.com \\
Website: https://sites.google.com/view/zhongli-wang \\
Bielefeld, Germany \\

\vspace{0.5cm}

% --- Recipient Information ---
\noindent
\textbf{Recruitment Committee} \\
Living Conditions Department \\
Luxembourg Institute of Socio-Economic Research (LISER) \\
Belval, Luxembourg \\

\vspace{0.5cm}

% --- Subject Line ---
\noindent
\textbf{Application for Research Associate (Postdoc) — Ref. 25-44} \\

\vspace{0.3cm}

\noindent Dear Members of the Recruitment Committee,\\[0.2cm]

I am writing to apply for the Research Associate (Postdoc) position in the Living Conditions Department (Ref. 25-44) at LISER. I am currently completing my Ph.D.\ in Economics at Bielefeld University and Università Cattolica del Sacro Cuore through the Marie Skłodowska–Curie EPOC programme. My research combines empirical microeconomics, causal machine learning, and computational social science to study inequality, technology adoption, and AI-mediated behavioral responses. The position at LISER aligns exceptionally well with my methodological background and my interest in understanding how digital transformation and AI systems affect economic behaviour, social mobility, and the distribution of opportunities.

Across my doctoral work, I have developed quantitative tools to analyse how individuals and firms adapt to new technologies, how digital infrastructures shape behavioural heterogeneity, and how policy interventions can reduce disparities in adoption and outcomes. My dissertation relies on a mixture of econometric analysis, machine-learning methods, and agent-based modelling to study technology-driven markets under heterogeneity and dynamic adjustment. These skills translate naturally to LISER’s research priorities, particularly those involving AI-enhanced econometrics, inequality measurement, and the analysis of behavioural microdata.

In addition to my dissertation, I have coauthored two projects that highlight my experience with machine learning and high-dimensional data.  
\emph{Polytope Fraud Theory} applies ML to detect anomalies in corporate financial records, demonstrating my ability to work with large datasets, develop predictive models, and ensure interpretability.  
A second project develops an agentic workflow that integrates large-language-model agents into empirical research, providing insight into human–AI interaction, automated text analysis, and scalable computational pipelines.  
These projects illustrate my ability to combine empirical rigour, algorithmic tools, and conceptual framing—an approach that fits LISER’s ambition to produce high-impact, policy-relevant research on emerging societal challenges.

I am strongly motivated to join LISER because of its excellence in empirical research, its interdisciplinary environment, and its commitment to studying inequality and social mobility using advanced data-driven methods. The Living Conditions Department, with its focus on distributional analysis, behavioural responses, and the societal impacts of technological change, provides an ideal environment to further develop my research agenda. I am particularly excited by the opportunity to contribute to work on:
\begin{itemize}
    \item the impact of AI and automation on inequality, discrimination, and equality of opportunity;
    \item behavioural adaptations in consumption, saving, and risk attitudes under technological change;
    \item micro-level dynamics of social and economic mobility using large-scale survey and administrative data;
    \item methodological advances in applied machine learning and AI-supported causal inference.
\end{itemize}

My experience working with large datasets, implementing reproducible computational workflows, and translating analytical results into clear narratives prepares me to contribute effectively to LISER’s research output, grant proposals, and policy-oriented publications.  
Because all chapters of my dissertation are single-authored, I am accustomed to independently designing and executing complex research projects while maintaining high methodological standards.

Thank you very much for considering my application. I would be delighted to contribute to LISER’s mission of producing rigorous, policy-relevant research on inequality and the social consequences of technological change, and I would welcome the opportunity to discuss my candidacy further.

\vspace{6mm}

\noindent Sincerely,\\[15mm]
\noindent Zhongli Wang

\end{document}
